% Options for packages loaded elsewhere
\PassOptionsToPackage{unicode}{hyperref}
\PassOptionsToPackage{hyphens}{url}
\PassOptionsToPackage{table}{xcolor}
%
\documentclass[
  11pt,
]{report}
\usepackage{amsmath,amssymb}
\usepackage{iftex}
\ifPDFTeX
  \usepackage[T1]{fontenc}
  \usepackage[utf8]{inputenc}
  \usepackage{textcomp} % provide euro and other symbols
\else % if luatex or xetex
  % pdfLaTeX-only repository policy: the original unicode-math/fontspec
  % branch is unreachable under pdfLaTeX and is replaced with a guard.
  \PackageError{cpython-suite}{Compile with pdfLaTeX only}{}
\fi
\usepackage{lmodern}
\ifPDFTeX\else
  % xetex/luatex font selection
\fi
% Use upquote if available, for straight quotes in verbatim environments
\IfFileExists{upquote.sty}{\usepackage{upquote}}{}
\IfFileExists{microtype.sty}{% use microtype if available
  \usepackage[]{microtype}
  \UseMicrotypeSet[protrusion]{basicmath} % disable protrusion for tt fonts
}{}
\makeatletter
\@ifundefined{KOMAClassName}{% if non-KOMA class
  \IfFileExists{parskip.sty}{%
    \usepackage{parskip}
  }{% else
    \setlength{\parindent}{0pt}
    \setlength{\parskip}{6pt plus 2pt minus 1pt}}
}{% if KOMA class
  \KOMAoptions{parskip=half}}
\makeatother
\usepackage{xcolor}
\usepackage{listings}
\newcommand{\passthrough}[1]{#1}
\lstset{defaultdialect=[5.3]Lua}
\lstset{defaultdialect=[x86masm]Assembler}
\usepackage{longtable,booktabs,array}
\usepackage{calc} % for calculating minipage widths
% Correct order of tables after \paragraph or \subparagraph
\usepackage{etoolbox}
\makeatletter
\patchcmd\longtable{\par}{\if@noskipsec\mbox{}\fi\par}{}{}
\makeatother
% Allow footnotes in longtable head/foot
\IfFileExists{footnotehyper.sty}{\usepackage{footnotehyper}}{\usepackage{footnote}}
\makesavenoteenv{longtable}
\setlength{\emergencystretch}{3em} % prevent overfull lines
\providecommand{\tightlist}{%
  \setlength{\itemsep}{0pt}\setlength{\parskip}{0pt}}
\setcounter{secnumdepth}{5}
% CPython Engineering Suite shared presentation layer -- pdfLaTeX compatible
\usepackage[T1]{fontenc}
\usepackage{lmodern}
\usepackage{microtype}
\usepackage{geometry}
\geometry{letterpaper,margin=0.91in,headheight=16pt,headsep=19pt}
\usepackage{xcolor} % [table] passed at top of file (option-clash fix)
\definecolor{CPNavy}{HTML}{13283B}
\definecolor{CPBlue}{HTML}{176BA0}
\definecolor{CPGold}{HTML}{D8A73C}
\definecolor{CPMuted}{HTML}{566979}
\definecolor{CPBg}{HTML}{EFF5F8}
\usepackage{titlesec}
\titleformat{\chapter}[display]{\normalfont\sffamily\bfseries\color{CPNavy}}{\large\color{CPBlue}CHAPTER \thechapter}{9pt}{\Huge}[\vspace{6pt}\color{CPGold}\titlerule]
\titleformat{\section}{\Large\bfseries\sffamily\color{CPBlue}}{\thesection}{0.7em}{}
\titleformat{\subsection}{\large\bfseries\sffamily\color{CPNavy}}{\thesubsection}{0.7em}{}
\usepackage{fancyhdr}
\pagestyle{fancy}
\fancyhf{}
\fancyhead[L]{\footnotesize\color{CPMuted}\sffamily CPYTHON ENGINEERING SUITE}
\fancyhead[R]{\footnotesize\color{CPMuted}\sffamily 3.16.0a1+dev SNAPSHOT}
\fancyfoot[C]{\footnotesize\color{CPMuted}\thepage}
\renewcommand{\headrulewidth}{0.3pt}
\AtBeginDocument{\urlstyle{tt}}
\lstset{basicstyle=\ttfamily\small,columns=fullflexible,keepspaces=true,breaklines=true,breakatwhitespace=false,showstringspaces=false,frame=single,framesep=5pt,rulecolor=\color{CPBg},backgroundcolor=\color{CPBg},xleftmargin=3pt,xrightmargin=3pt}
\usepackage{tcolorbox}
\tcbuselibrary{breakable,skins}
\newtcolorbox{sourcebox}[1][]{enhanced,breakable,sharp corners,boxrule=0.4pt,colback=CPBg,colframe=CPBlue,borderline west={3pt}{0pt}{CPBlue},left=10pt,right=8pt,top=6pt,bottom=6pt,#1}
\newtcolorbox{cautionbox}[1][]{enhanced,breakable,sharp corners,boxrule=0.4pt,colback=CPBg,colframe=CPGold,borderline west={3pt}{0pt}{CPGold},left=10pt,right=8pt,top=6pt,bottom=6pt,#1}
\usepackage{tikz}
\usetikzlibrary{arrows.meta,positioning,shapes.geometric,calc}
\usepackage{enumitem}
\setlist{itemsep=2pt,topsep=4pt}
\usepackage{booktabs}
\usepackage{longtable}
\AtBeginEnvironment{longtable}{\small}
\usepackage{array}
\setlength{\parskip}{5pt}
\setlength{\parindent}{0pt}
\setcounter{tocdepth}{2}
\usepackage{hyperref}
\hypersetup{colorlinks=true,linkcolor=CPBlue,urlcolor=CPBlue,citecolor=CPBlue,pdfauthor={CPython Engineering Documentation Suite}}
\makeatletter
\renewcommand{\maketitle}{%
  \begin{titlepage}\pagecolor{CPNavy}\color{white}\sffamily
  \vspace*{0.8in}
  {\color{CPGold}\rule{2.1in}{4pt}}\\[16pt]
  {\Large\bfseries CPYTHON ENGINEERING DOCUMENTATION SUITE}\\[24pt]
  {\Huge\bfseries \@title\par}\vspace{19pt}
  \vfill
  {\normalsize Source baseline: CPython 3.16.0 alpha 1+dev repository snapshot}\par
  {\normalsize Prepared 16 October 2026}\par
  {\small\color{white!70} Source-traced engineering commentary; not an official Python specification.}\vspace{0.65in}
  \end{titlepage}\nopagecolor\color{black}
  }
\makeatother
\ifLuaTeX
  \usepackage{selnolig}  % disable illegal ligatures
\fi
\usepackage{bookmark}
\IfFileExists{xurl.sty}{\usepackage{xurl}}{} % add URL line breaks if available
\urlstyle{same}
\hypersetup{
  pdftitle={Python Standard Library Engineering Guide},
  hidelinks,
  pdfcreator={LaTeX via pandoc}}

\title{Python Standard Library Engineering Guide}
\usepackage{etoolbox}
\makeatletter
\providecommand{\subtitle}[1]{% add subtitle to \maketitle
  \apptocmd{\@title}{\par {\large #1 \par}}{}{}
}
\makeatother
\subtitle{Volume 02 \textbar{} Core APIs, data pipelines, concurrency
and reliable services}
\author{}
\date{}

\begin{document}
\maketitle

{
\setcounter{tocdepth}{1}
\tableofcontents
}
\chapter{Scope and module-selection
discipline}\label{scope-and-module-selection-discipline}

\begin{sourcebox}
\textbf{Evidence classification.} \textbf{SOURCE} = statements supported by material present in the uploaded snapshot; \textbf{IMPLEMENTATION} = CPython-specific internals, not guaranteed Python semantics; \textbf{ENGINEERING} = original recommendations, examples, checks, and exercises. The source baseline is \textbf{3.16.0 alpha 1+dev}; missing repository file bodies are noted as evidence limitations.
\end{sourcebox}

Python's standard library exposes interfaces spanning data structures,
I/O, system integration, networking, formats, process control,
concurrency, testing and cryptographic primitives. In the uploaded
snapshot, \passthrough{\lstinline!Doc/library/!} provides
module-specific contracts and \passthrough{\lstinline!Lib/!} supplies
many reference implementations; some C accelerators reside under
\passthrough{\lstinline!Modules/!}.
\path{InternalDocs/structure.md} explicitly describes
this familiar documentation, source, test and native-accelerator layout.

\textbf{SOURCE} means the behavior is grounded in the corresponding
module's documentation. \textbf{ENGINEERING} means a selection, pattern
or test proposed here. Avoid treating deprecated or removed modules as
recommended dependencies simply because legacy documentation or
compatibility references still exist in a development branch.

\begin{longtable}[]{@{}
  >{\raggedright\arraybackslash}p{(\columnwidth - 4\tabcolsep) * \real{0.3333}}
  >{\raggedright\arraybackslash}p{(\columnwidth - 4\tabcolsep) * \real{0.3333}}
  >{\raggedright\arraybackslash}p{(\columnwidth - 4\tabcolsep) * \real{0.3333}}@{}}
\toprule\noalign{}
\begin{minipage}[b]{\linewidth}\raggedright
Capability
\end{minipage} & \begin{minipage}[b]{\linewidth}\raggedright
First modules to assess
\end{minipage} & \begin{minipage}[b]{\linewidth}\raggedright
Design concern
\end{minipage} \\
\midrule\noalign{}
\endhead
\bottomrule\noalign{}
\endlastfoot
Paths and filesystem & \passthrough{\lstinline!pathlib!},
\passthrough{\lstinline!os!}, \passthrough{\lstinline!shutil!},
\passthrough{\lstinline!tempfile!} & Encoding, traversal, lifetime,
permissions \\
Data collections & \passthrough{\lstinline!collections!},
\passthrough{\lstinline!heapq!}, \passthrough{\lstinline!bisect!},
\passthrough{\lstinline!itertools!} & Algorithmic complexity and
mutation \\
Structured records & \passthrough{\lstinline!dataclasses!},
\passthrough{\lstinline!enum!}, \passthrough{\lstinline!typing!} &
Runtime versus static contracts \\
Text and formats & \passthrough{\lstinline!re!},
\passthrough{\lstinline!csv!}, \passthrough{\lstinline!json!},
\passthrough{\lstinline!tomllib!} & Parsing, memory and schema
validation \\
Time and arithmetic & \passthrough{\lstinline!datetime!},
\passthrough{\lstinline!zoneinfo!}, \passthrough{\lstinline!decimal!},
\passthrough{\lstinline!fractions!} & Representation, timezone,
rounding \\
Network/I/O & \passthrough{\lstinline!socket!},
\passthrough{\lstinline!ssl!}, \passthrough{\lstinline!http.client!},
\passthrough{\lstinline!asyncio!} & TLS, timeouts, stream
backpressure \\
Parallel work & \passthrough{\lstinline!threading!},
\path{concurrent.futures},
\passthrough{\lstinline!multiprocessing!} & Isolation, thread safety,
IPC \\
Security primitives & \passthrough{\lstinline!secrets!},
\passthrough{\lstinline!hashlib!}, \passthrough{\lstinline!hmac!} &
Threat model, key lifecycle, misuse \\
Diagnostics & \passthrough{\lstinline!logging!},
\passthrough{\lstinline!tracemalloc!},
\passthrough{\lstinline!cProfile!} & Observation overhead and sensitive
data \\
Testing & \passthrough{\lstinline!unittest!},
\path{unittest.mock},
\passthrough{\lstinline!doctest!} & Isolation, reproducibility \\
\end{longtable}

\chapter{Files, encodings and transactional
boundaries}\label{files-encodings-and-transactional-boundaries}

\section{The path/object boundary}\label{the-pathobject-boundary}

Prefer \passthrough{\lstinline!pathlib.Path!} for high-level path
operations but do not confuse path construction with authorization.
Normalize your application's concept of an allowed root and scrutinize
links and race conditions when writing files influenced by an adversary.
\path{Path.resolve()} alone does not remove all
time-of-check/time-of-use risks. External text should specify encoding;
binary data should not be silently decoded using process defaults.

\begin{lstlisting}[language=Python]
from pathlib import Path

def load_config(path: Path) -> str:
    if not path.is_file():
        raise FileNotFoundError(path)
    return path.read_text(encoding="utf-8")
\end{lstlisting}

\section{Atomic replacement for completed
outputs}\label{atomic-replacement-for-completed-outputs}

\textbf{ENGINEERING example:} create a file in the target directory,
flush and close it, then replace the destination. Replace semantics and
persistence guarantees vary by platform and filesystem; durability
across power loss may require directory syncing and other
platform-specific work.

\begin{lstlisting}[language=Python]
import os
from pathlib import Path
from tempfile import NamedTemporaryFile

def replace_text(destination: Path, content: str) -> None:
    destination.parent.mkdir(parents=True, exist_ok=True)
    temp = None
    try:
        with NamedTemporaryFile(mode="w", encoding="utf-8",
                                dir=destination.parent, delete=False) as fp:
            temp = Path(fp.name)
            fp.write(content)
            fp.flush()
            os.fsync(fp.fileno())
        os.replace(temp, destination)
    finally:
        if temp is not None:
            temp.unlink(missing_ok=True)
\end{lstlisting}

\chapter{Data structures, complexity and practical
contracts}\label{data-structures-complexity-and-practical-contracts}

\passthrough{\lstinline!dict!}, \passthrough{\lstinline!list!},
\passthrough{\lstinline!set!}, \passthrough{\lstinline!tuple!}, and
\passthrough{\lstinline!deque!} serve different operational goals.
\path{Doc/builtins/time-complexity.rst} is a source
of complexity expectations in the snapshot; complexity descriptions are
not exact per-operation latency guarantees. A
\passthrough{\lstinline!deque!} is useful for two-ended queue
operations; a list is useful for indexed random access. Sorting can use
a \passthrough{\lstinline!key!} function rather than mutating records or
implementing a custom comparator.

\begin{lstlisting}[language=Python]
from collections import Counter, deque

queue = deque(["validate", "build"])
queue.append("test")
assert queue.popleft() == "validate"
counts = Counter("abac")
assert counts["a"] == 2
\end{lstlisting}

\textbf{ENGINEERING checklist:} measure typical and worst-case sizes,
make ordering requirements explicit, verify behavior under concurrent
mutation, and avoid using lists as unbounded message queues without an
admission policy.

\chapter{Parsing, serialization and trust
boundaries}\label{parsing-serialization-and-trust-boundaries}

\section{JSON, TOML, CSV and pickle}\label{json-toml-csv-and-pickle}

\passthrough{\lstinline!json!} is suitable for interoperable data
exchange, but is not a schema validator.
\passthrough{\lstinline!tomllib!} loads TOML; it is not a
standard-library TOML writer. \passthrough{\lstinline!csv!} parsing
requires appropriate newline treatment and application-level limits for
field length and row count. \passthrough{\lstinline!pickle!} is a
Python-object serialization mechanism that \textbf{must not be used to
deserialize untrusted input}: its own documentation warns that
unpickling malicious data can execute code.

\begin{lstlisting}[language=Python]
import json

def parse_request(raw: str) -> dict:
    payload = json.loads(raw)
    if not isinstance(payload, dict):
        raise ValueError("object expected")
    name = payload.get("name")
    if not isinstance(name, str) or not (1 <= len(name) <= 120):
        raise ValueError("invalid name")
    return {"name": name}
\end{lstlisting}

Validation can require stronger constraints (size ceilings,
duplicate-key policy, Unicode normalization, numeric range, structure
depth) than the parser itself supplies. Structured input must never
become trusted code solely because it passed syntactic decoding.

\chapter{Concurrency design and
scheduling}\label{concurrency-design-and-scheduling}

\section{Distinguish concurrency
mechanisms}\label{distinguish-concurrency-mechanisms}

\passthrough{\lstinline!asyncio!} is documented as a library for
concurrent
\passthrough{\lstinline!async!}/\passthrough{\lstinline!await!} programs
and is particularly suited to high-level I/O-bound work. Threads can
interleave operations, and in a free-threaded CPython build can execute
Python work in parallel; synchronization is still required for
correctness. Processes isolate runtime memory but introduce IPC cost and
serialization considerations.
\path{concurrent.futures} provides an abstraction
over thread and process execution.

\begin{lstlisting}[language=Python]
import asyncio

async def fetch_with_timeout(operation, seconds: float):
    async with asyncio.timeout(seconds):
        return await operation()

async def run_jobs(jobs):
    async with asyncio.TaskGroup() as group:
        tasks = [group.create_task(job()) for job in jobs]
    return [task.result() for task in tasks]
\end{lstlisting}

The sample assumes \passthrough{\lstinline!jobs!} contains coroutine
factories and targets runtimes supporting the shown asyncio APIs; use
explicit cancellation and timeout tests.
\passthrough{\lstinline!TaskGroup!} supports structured lifetime
management: errors from child tasks propagate according to group
semantics, which may involve exception groups.

\section{Thread/process decision
record}\label{threadprocess-decision-record}

\begin{longtable}[]{@{}
  >{\raggedright\arraybackslash}p{(\columnwidth - 6\tabcolsep) * \real{0.2500}}
  >{\raggedright\arraybackslash}p{(\columnwidth - 6\tabcolsep) * \real{0.2500}}
  >{\raggedright\arraybackslash}p{(\columnwidth - 6\tabcolsep) * \real{0.2500}}
  >{\raggedright\arraybackslash}p{(\columnwidth - 6\tabcolsep) * \real{0.2500}}@{}}
\toprule\noalign{}
\begin{minipage}[b]{\linewidth}\raggedright
Question
\end{minipage} & \begin{minipage}[b]{\linewidth}\raggedright
Async tasks
\end{minipage} & \begin{minipage}[b]{\linewidth}\raggedright
Threads
\end{minipage} & \begin{minipage}[b]{\linewidth}\raggedright
Processes
\end{minipage} \\
\midrule\noalign{}
\endhead
\bottomrule\noalign{}
\endlastfoot
Waiting on large numbers of sockets? & Often effective & Sometimes
useful & Usually high overhead \\
CPU-bound work? & Not inherently parallel & Build/extension-dependent &
Common default isolation path \\
Shared mutable state? & Still needs disciplined access &
Lock/synchronization design required & State copied or passed through
IPC \\
Cancellation? & Cooperative task cancellation & Application-specific
coordination & Termination/IPC strategy \\
Main risk & Blocking event loop & Races and deadlocks & Serialization
and process lifecycle \\
\end{longtable}

\chapter{Time, numeric integrity and
secrets}\label{time-numeric-integrity-and-secrets}

Wall-clock time, monotonic elapsed time, timezone-aware date-time
values, and durations answer different questions. Prefer
\path{time.monotonic()} for elapsed-time thresholds
and explicit \passthrough{\lstinline!zoneinfo!} timezones for civil
time. Use \path{decimal.Decimal} where base-ten
precision and rounding policy are contractual; construction from binary
floating-point may import an already-rounded value.

\begin{lstlisting}[language=Python]
from decimal import Decimal, ROUND_HALF_UP

def invoice_total(items: list[tuple[str, str]]) -> Decimal:
    total = sum((Decimal(amount) for _, amount in items), Decimal("0"))
    return total.quantize(Decimal("0.01"), rounding=ROUND_HALF_UP)
\end{lstlisting}

\passthrough{\lstinline!random!} is for simulations and nonsecurity
applications, \textbf{not} generating credentials. The snapshot's
\path{Doc/library/security_warnings.rst} directs
security uses to \passthrough{\lstinline!secrets!}.
\passthrough{\lstinline!hashlib!} and \passthrough{\lstinline!hmac!} are
primitives, not a complete password storage design; use a suitable
dedicated password-hashing scheme in an actual authentication system.

\begin{lstlisting}[language=Python]
import secrets

reset_token = secrets.token_urlsafe(32)
assert isinstance(reset_token, str)
\end{lstlisting}

\chapter{Resource boundaries, observability and operational
APIs}\label{resource-boundaries-observability-and-operational-apis}

\section{Logging without data
exposure}\label{logging-without-data-exposure}

Use named loggers, parameterized log messages, coherent event context
and structured filters. The docs warn that some logging configuration
mechanisms may evaluate expressions; treat logging configuration as
trusted code or carefully controlled input. Never log passwords, auth
tokens, high-sensitivity payloads or raw personal data unless an
explicit retention and authorization policy requires it.

\begin{lstlisting}[language=Python]
import logging
logger = logging.getLogger(__name__)

def process_batch(batch_id: str, count: int) -> None:
    logger.info("batch accepted id=%s records=%d", batch_id, count)
\end{lstlisting}

\section{Network hygiene}\label{network-hygiene}

Treat \passthrough{\lstinline!http.server!} as a teaching utility rather
than production-hardened HTTP serving software (the library docs call
out limited security checks). Set explicit network timeouts, make TLS
hostname/certificate checking part of the client contract, and avoid
exposing debug endpoints. Reject weak fallback behavior rather than
silently turning off TLS verification when a connection fails.

\chapter{Deployable reference
mini-project}\label{deployable-reference-mini-project}

\textbf{Project:} a log-safe, configuration-driven, asynchronous batch
ingester.

\begin{enumerate}
\def\labelenumi{\arabic{enumi}.}
\tightlist
\item
  Accept JSON input with a fixed schema and size ceiling; reject
  unexpected types.
\item
  Load TOML application configuration with
  \passthrough{\lstinline!tomllib!} and validate it.
\item
  Queue bounded units of work and use
  \path{asyncio.TaskGroup} for managed concurrency.
\item
  Wrap each external I/O in an explicit timeout and propagate
  cancellation.
\item
  Save completed output using temporary-file replacement, with an
  documented durability policy.
\item
  Log a generated request correlation ID but not raw secrets or
  sensitive payloads.
\item
  Test parse failures, file-permission failures, timeouts, cancellation,
  partial writes and restart behavior.
\end{enumerate}

\textbf{Acceptance:} automated tests exercise positive, negative and
concurrency paths; operator guide specifies file permissions, timeout
defaults, maximum request size and secret-handling expectations.

\chapter{Source concordance}\label{source-concordance}

\begin{itemize}
\tightlist
\item
  \textbf{SOURCE:} \path{InternalDocs/structure.md};
  \path{Doc/library/index.rst},
  \path{Doc/library/asyncio.rst},
  \path{Doc/library/venv.rst},
  \path{Doc/library/pickle.rst},
  \path{Doc/library/security_warnings.rst}.
\item
  \textbf{SOURCE:} \path{Doc/library/json.rst},
  \passthrough{\lstinline!tomllib.rst!},
  \passthrough{\lstinline!csv.rst!},
  \passthrough{\lstinline!pathlib.rst!} (where present in source
  listing), \passthrough{\lstinline!secrets.rst!},
  \passthrough{\lstinline!logging.rst!},
  \path{concurrent.futures.rst},
  \passthrough{\lstinline!datetime.rst!}.
\item
  \textbf{SOURCE:} \passthrough{\lstinline!Lib/json/!},
  \passthrough{\lstinline!Lib/asyncio/!},
  \passthrough{\lstinline!Lib/venv/!},
  \passthrough{\lstinline!Lib/test/!} when included in the text export.
\item
  \textbf{ENGINEERING:} examples, tradeoff table, operational
  guardrails, and mini-project are synthesized for this guide, not
  guarantees expressed by every module.
\end{itemize}

\textbf{Version discipline:} cross-check available modules against the
specific interpreter target and its
\path{Doc/deprecations/} inventory; the presence of
historical documentation is not evidence a module remains supported.

\end{document}
